\documentclass[11pt,a4paper]{article}
\usepackage[utf8]{inputenc}
\usepackage[margin=2.5cm]{geometry}
\usepackage{amsmath,amssymb}
\usepackage{booktabs}
\usepackage{graphicx}
\usepackage{hyperref}
\usepackage{listings}
\usepackage{xcolor}
\usepackage{float}
\usepackage{siunitx}

\lstset{
  basicstyle=\ttfamily\small,
  keywordstyle=\color{blue},
  commentstyle=\color{gray},
  stringstyle=\color{red},
  breaklines=true,
  frame=single,
  numbers=left,
  numberstyle=\tiny\color{gray}
}

\title{MPC Optimization Findings\\
\large F1Tenth Bachelor Project --- CPU, FPGA, and Simulator}
\author{Optimization Report}
\date{\today}

\begin{document}
\maketitle

\tableofcontents
\newpage

%% ============================================================
\section{Overview}
This document summarizes the optimizations and findings for the Model Predictive Control (MPC) implementation across three targets: CPU (Task~B), FPGA (Task~A), and the F1Tenth simulator (Task~C). The MPC uses a Riccati-ADMM solver with a 20-step prediction horizon and \SI{0.04}{\second} time step (\SI{0.8}{\second} lookahead) on a Frenet-frame single-track vehicle model. The FPGA implementation uses 19 steps (\SI{0.76}{\second} lookahead) due to resource constraints.

\subsection{Vehicle Parameters}
All implementations use the same measured vehicle parameters (Table~\ref{tab:vehicle_params}).

\begin{table}[H]
\centering
\caption{Measured vehicle parameters used across all implementations.}
\label{tab:vehicle_params}
\begin{tabular}{@{}llrl@{}}
\toprule
\textbf{Parameter} & \textbf{Symbol} & \textbf{Value} & \textbf{Source} \\
\midrule
Mass & $m$ & \SI{3.314}{\kilo\gram} & Measured \\
Moment of inertia & $I_z$ & \SI{0.035}{\kilo\gram\meter\squared} & CAD \\
Front axle distance & $l_f$ & \SI{0.166}{\meter} & CAD \\
Rear axle distance & $l_r$ & \SI{0.160}{\meter} & CAD \\
Wheelbase & $L$ & \SI{0.326}{\meter} & $l_f + l_r$ \\
CG height & $h$ & \SI{0.0703}{\meter} & CAD \\
Friction coefficient & $\mu$ & 0.7463 & Tested \\
Front tire stiffness & $C_{S_f}$ & 2.804 & Tested \\
Rear tire stiffness & $C_{S_r}$ & 3.320 & Tested \\
Max steering angle & $\delta_{\max}$ & \SI{0.4282}{\radian} & Tested \\
Max steering velocity & $\dot{\delta}_{\max}$ & \SI{2.8492}{\radian\per\second} & Tested \\
Max acceleration & $a_{\max}$ & \SI{8.0}{\meter\per\second\squared} & Tested \\
Power transition speed & $v_{\text{switch}}$ & \SI{7.319}{\meter\per\second} & Default \\
\bottomrule
\end{tabular}
\end{table}

%% ============================================================
\section{Task B: CPU MPC Optimization}

The CPU MPC targets the Jetson Xavier NX (ARM Cortex-A57 @ \SI{1.9}{\giga\hertz}), tested on an x86 laptop.

\subsection{Float32 Conversion}
\label{sec:float32}
The original CPU implementation used Q16.16 fixed-point arithmetic (matching the FPGA). This was converted to native \texttt{float32} operations:

\begin{itemize}
  \item \texttt{typedef float fixed\_point\_t}. All \texttt{fp\_*} functions became inline float operations.
  \item Removed Taylor series and Newton-Raphson approximations (replaced by native FPU ops).
  \item All \texttt{int64\_t} intermediates and \texttt{>> FP\_FRAC\_BITS} shifts eliminated.
\end{itemize}

\subsection{Speed Progression}

\begin{table}[H]
\centering
\caption{CPU solve time progression.}
\label{tab:cpu_speed}
\begin{tabular}{@{}lrrr@{}}
\toprule
\textbf{Configuration} & \textbf{Avg Time} & \textbf{Reduction} & \textbf{Tests} \\
\midrule
Baseline (Q16.16, -O2)      & \SI{111.7}{\micro\second} & ---    & 6/6 PASS \\
Float32, -O2                 & \SI{64.0}{\micro\second}  & 43\%  & 6/6 PASS \\
Float32, -O3 -ffast-math     & \SI{39.0}{\micro\second}  & 65\%  & 6/6 PASS \\
+ Rho persistence            & \SI{11.8}{\micro\second}  & 89\%  & 6/6 PASS \\
+ Tolerance 5.0, warm-start  & \SI{5.3}{\micro\second}   & 95\%  & 6/6 PASS \\
\bottomrule
\end{tabular}
\end{table}

\subsection{Key Speed Optimizations}
\begin{enumerate}
  \item \textbf{Persist ADMM~$\rho$ across warm-starts.} The CPU previously reset $\rho = 25.0$ every call. After fixing this, average iterations dropped from 11.0 to 3.2 ($-71\%$). Single largest improvement.
  \item \textbf{Direct Frenet linearization.} Eliminated global-to-Frenet indirection (full $6\times 6$ matrices $\to$ direct Frenet-frame linearization).
  \item \textbf{Sparse zeroing.} Replaced \texttt{memset(step\_data, 0, 31\,KB)} with targeted field zeroing.
  \item \textbf{Tolerance relaxation.} With warm-start quality, tolerance 5.0 yields 1.0 average iterations at \SI{200}{\hertz}.
\end{enumerate}

\subsection{MPC Control Frequency}
200\,Hz is optimal. Higher rates (400--500\,Hz) provide no measurable improvement in tracking accuracy.

%% ============================================================
\section{Task A: FPGA MPC Optimization}

The FPGA targets the Ultra96-V2 (Xilinx ZU3EG) at \SI{100}{\mega\hertz} using Q16.16 fixed-point arithmetic.

\subsection{Algorithmic Optimizations}
Applied in order:
\begin{enumerate}
  \item \textbf{B-matrix sparsity}: Exploiting sparse structure of the Frenet linearized B matrix to skip zero multiplications.
  \item \textbf{Persistent rho}: Maintaining ADMM penalty parameter $\rho$ across solver calls (warm-starting the dual).
  \item \textbf{Adaptive rho every 2 iterations}: Reduced check frequency.
  \item \textbf{Over-relaxation simplification}: $\hat{x} = \alpha x + (1-\alpha)z \to \hat{x} = x + (\alpha - 1)(x - z)$. Saves 208~DSP multiply operations per ADMM iteration.
\end{enumerate}

\subsection{HLS Pipeline Optimizations}
\begin{enumerate}
  \item \textbf{\texttt{fp\_mul} non-inline 4-cycle DSP binding}: Replaced inline multiplication with a dedicated 4-cycle pipelined DSP48E2 multiplier (\texttt{BIND\_OP impl=dsp latency=3}), controlled by \texttt{MPC\_HLS\_MUL\_LIMIT=4}. This prevents DSP chaining and resolves critical-path timing violations.
  \item \textbf{\texttt{fp\_normalize\_angle} DSP fix}: Replaced branch-based normalization with arithmetic-only implementation to avoid DSP-to-logic critical paths.
  \item \textbf{Step~5 II=2$\to$1}: The Riccati backward pass Step~5 (P-matrix update) achieved initiation interval~1, doubling throughput for this step (only 2~multiplies, within MUL\_LIMIT=4).
  \item \textbf{P~symmetry exploitation}: Reduced Step~5 from $8\times8=64$ iterations to 21 upper-triangle iterations with static decode tables and mirror writes ($P[j][i] = P[i][j]$).
  \item \textbf{G~sparsity for P~rows~6,7}: Exploiting that only 2 of 16 entries in the N\_cross matrix are non-zero (G[0][6] and G[1][7]), reducing multiplications by half and enabling II=1 for these rows.
\end{enumerate}

\subsection{Accuracy Fixes}
\begin{enumerate}
  \item \textbf{Reciprocal CLZ seed bug (critical)}: The initial Newton-Raphson seed for \texttt{fp\_recip} used \texttt{\_\_builtin\_clz(abs\_x) - 1} instead of \texttt{\_\_builtin\_clz(abs\_x)}. Similarly, \texttt{reciprocal\_64} used \texttt{\_\_builtin\_clzll - 33} instead of \texttt{-32}. This caused the initial guess to be $4\times$ too low, resulting in 5--10\% systematic error after 3~NR iterations. After fix: $<0.002\%$ error ($5000\times$ improvement). NR iterations restored to 4 (from 3, which was insufficient with the wrong seed).
  \item \textbf{Forward pass single-shift accumulation}: Inner products in the forward pass previously shifted each product individually before accumulation ($\sim$6~LSB loss per 6-term sum). Changed to accumulate full \texttt{int64\_t} products with a single shift at the end. Worst-case overflow: $\sim 2^{47}$, safely within \texttt{int64\_t} range.
\end{enumerate}

\textbf{Impact:} Average lateral tracking error improved 19\% (0.144$\to$\SI{0.117}{\metre}) in FPGA simulation.

\subsection{FPGA HLS Synthesis Results}
Synthesized with Vitis HLS~2025.1 targeting ZU3EG at \SI{100}{\mega\hertz}.

\begin{table}[H]
\centering
\caption{FPGA resource utilization (HLS estimates and Vivado implementation).}
\label{tab:fpga_synthesis}
\begin{tabular}{@{}lrrrr@{}}
\toprule
\textbf{Resource} & \textbf{HLS Est.} & \textbf{Vivado Impl.} & \textbf{Available} & \textbf{Vivado Util\%} \\
\midrule
DSP48E2   & 346  & 321    & 360     & 89.2\% \\
LUT       & 60,221 & 46,934 & 70,560 & 66.5\% \\
FF        & ---  & 39,081 & 141,120 & 27.7\% \\
CLB       & ---  & 8,491  & 8,820   & 96.3\% \\
BRAM (18K)& 90   & ---    & 432     & 21\% \\
\bottomrule
\end{tabular}
\end{table}

\begin{table}[H]
\centering
\caption{FPGA timing and latency.}
\label{tab:fpga_timing}
\begin{tabular}{@{}lr@{}}
\toprule
\textbf{Metric} & \textbf{Value} \\
\midrule
Target clock period         & \SI{10}{\nano\second} (\SI{100}{\mega\hertz}) \\
HLS estimated $F_\text{max}$    & \SI{113}{\mega\hertz} \\
Vivado WNS (post-route)    & \SI{+0.044}{\nano\second} (timing closed) \\
Worst-case latency (8 ADMM iters) & 82,786 cycles = \SI{0.828}{\milli\second} \\
Typical latency (1.1 avg iters)   & $\sim$12,000 cycles $\approx$ \SI{0.12}{\milli\second} \\
Per-iteration latency       & $\sim$11,200 cycles $\approx$ \SI{0.112}{\milli\second} \\
Control budget (\SI{200}{\hertz}) & \SI{5}{\milli\second} (83\% headroom) \\
\bottomrule
\end{tabular}
\end{table}

\subsection{Synthesis Strategy}
Vivado synthesis uses \texttt{Flow\_PerfOptimized\_high}; implementation uses \texttt{Performance\_ExtraTimingOpt}. The \texttt{MPC\_HLS\_MUL\_LIMIT=4} controls DSP multiplier instances. Testing with MUL\_LIMIT=5 produced \emph{worse} timing due to increased routing congestion.

%% ============================================================
\section{Task C: Simulator Accuracy}

\subsection{Dynamics Model Alignment}
The simulation testbench matches the f1tenth\_gym's single-track (ST) dynamic model exactly:
\begin{itemize}
  \item \textbf{Model}: 7-state single-track $[X, Y, \delta, V, \psi, \dot{\psi}, \beta]$
  \item \textbf{Integrator}: RK4, matching gym's integrator.py
  \item \textbf{Steering}: Bang-bang PID controller converting desired angle to steering velocity at $\pm\dot{\delta}_{\max}$
  \item \textbf{Kinematic/Dynamic switch}: At $V < \SI{0.5}{\meter\per\second}$, kinematic bicycle; above, full tire-force dynamics
  \item \textbf{Slip angles}: \texttt{atan2}-based (nonlinear), matching gym exactly
  \item \textbf{Force resolution}: $\cos(\delta)/\sin(\delta)$ terms for front tire forces
\end{itemize}

\subsection{Raceline Pipeline}
\label{sec:raceline_pipeline}
A unified script \texttt{optimize\_trajectory.py} provides one-command trajectory generation:
\[
  \text{Map (PGM)} \;\to\; \text{Boundary extraction} \;\to\; \text{Centerline (periodic B-spline)} \;\to\; \text{Width (ray-cast)} \;\to\; \text{TUM min-curvature optimizer} \;\to\; \text{MPC CSV}
\]
Key features: works on both designed maps and SLAM maps, auto-scales point count based on perimeter, curvature-based width capping ($w \leq 0.8\,R$), heading convention $\psi_{\text{TUM}} + \pi/2 = \psi_{\text{MPC}}$.

\subsection{Decoupled Simulation and Control Frequencies}
Independent sim physics and MPC control rates via environment variables. 1000\,Hz sim with 200\,Hz MPC recommended for modelling real-world continuous dynamics. All configurations pass 6/6 tests.

%% ============================================================
\section{Realistic Simulation Mode}
\label{sec:realistic_sim}

The \texttt{REALISTIC\_SIM=1} flag activates four physical effects simultaneously:
\begin{enumerate}
  \item \textbf{Drivetrain model} --- aerodynamic drag ($C_d A = 0.008$), rolling resistance ($\mu_r = 0.015$), and motor torque limits.
  \item \textbf{Pacejka tire model} --- non-linear $F_y(\alpha)$ with shape parameter $C = 1.9$ and $\mu = 0.7463$.
  \item \textbf{One-step actuator delay} --- 2-sample steering buffer delaying commanded angle by one MPC step (\SI{5}{\milli\second}).
  \item \textbf{Sensor noise} --- Gaussian noise on lateral error ($\sigma_e = \SI{3}{\milli\meter}$), heading ($\sigma_\psi = 0.004\;\text{rad}$), and velocity ($\sigma_v = \SI{10}{\milli\meter\per\second}$), with EMA low-pass filter ($\alpha = 0.35$).
\end{enumerate}

Maximum lateral acceleration updated to friction-limited value: $a_{\text{lat,max}} = \mu \cdot g = \SI{7.3212}{\metre\per\second\squared}$.

%% ============================================================
\section{Wall Constraint System}
\label{sec:wall_constraints}

The ADMM solver supports both hard box constraints and soft (quadratic penalty) wall constraints. Hard constraints ($k = 0$) use direct projection in the z-update; soft constraints ($k > 0$) use the proximal operator:
\[
  \text{prox}_{g/\rho}(v) =
  \begin{cases}
    \dfrac{k \cdot \text{bound} + \rho \cdot v}{k + \rho} & \text{if } v \text{ violates bound,} \\[4pt]
    v & \text{otherwise.}
  \end{cases}
\]

Current best configuration uses \textbf{hard constraints} (WALL\_SOFT\_K=0) with WALL\_MARGIN=\SI{0.15}{\metre} and WALL\_END=18, constraining the first 18 horizon steps.

%% ============================================================
\section{Weight Tuning}
\label{sec:weight_tuning}

Automated sweeps tested $>$1800 configurations under \texttt{REALISTIC\_SIM=1} across raceline variants (cl020, cl030, cl045, cl050).

\subsection{Current Best Configuration}
\begin{table}[H]
\centering
\caption{Best MPC weights (cl050 raceline, Spielberg, 100\,s, \SI{200}{\hertz}).}
\label{tab:best_weights}
\begin{tabular}{@{}lrl@{}}
\toprule
\textbf{Parameter} & \textbf{Value} & \textbf{Notes} \\
\midrule
$Q_\text{lat}$           & 340    & Lateral error weight \\
$Q_\text{hdg}$           & 1000   & Heading error weight \\
$Q_\text{vel}$           & 26     & Velocity tracking weight \\
$Q_\text{lat\_vel}$      & 69     & Lateral velocity weight \\
$Q_\text{yaw}$           & 22     & Yaw rate weight \\
$R_\text{steer}$         & 0.15   & Steering effort penalty \\
$R_\text{accel}$         & 0.01   & Acceleration effort penalty \\
$W_\text{jerk}$          & 0.3    & Steering rate penalty \\
$W_\text{accel\_rate}$   & 0.1    & Acceleration rate penalty \\
Prediction dt             & \SI{0.04}{\second} & 20 steps $\times$ 0.04\,s = 0.8\,s \\
Wall margin              & \SI{0.15}{\metre} & Hard box constraints \\
Wall end                 & 18     & Constrain first 18 steps \\
Wall soft stiffness      & 0      & Hard constraints (no penalty) \\
\bottomrule
\end{tabular}
\end{table}

\subsection{Raceline Variants}
\begin{table}[H]
\centering
\caption{Raceline variants (Spielberg track, $v_\text{max} = \SI{12}{\metre\per\second}$).}
\label{tab:raceline_variants}
\begin{tabular}{@{}lcccl@{}}
\toprule
\textbf{Variant} & \textbf{Clearance} & \textbf{Waypoints} & \textbf{MPC Tests} & \textbf{Notes} \\
\midrule
Tight (cl020)        & \SI{0.20}{\metre} & 966 & 5/6 & Low speed \\
Medium (cl030)       & \SI{0.30}{\metre} & 967 & 5/6 & Low speed \\
Default (cl045)      & \SI{0.45}{\metre} & 969 & 6/6 & Good balance \\
Wide (cl050)         & \SI{0.50}{\metre} & 970 & 6/6 & Highest speed, primary test target \\
\bottomrule
\end{tabular}
\end{table}

%% ============================================================
\section{Final System Performance}
\label{sec:final_performance}

\begin{table}[H]
\centering
\caption{Final system performance (cl050 raceline, Spielberg, 100\,s, \SI{200}{\hertz}).}
\label{tab:final_perf}
\begin{tabular}{@{}lrr@{}}
\toprule
\textbf{Metric} & \textbf{CPU MPC} & \textbf{FPGA MPC} \\
\midrule
Avg solve time          & \SI{5.3}{\micro\second}  & \SI{828}{\micro\second} (worst-case) \\
                        &                          & \SI{120}{\micro\second} (typical) \\
Avg iterations          & 1.0                      & 1.1 \\
Max iterations          & 6                        & 6 \\
Max velocity            & \SI{11.91}{\metre\per\second} & \SI{12.00}{\metre\per\second} \\
Avg lateral error       & \SI{0.110}{\metre}       & \SI{0.261}{\metre} \\
Max lateral error       & \SI{0.452}{\metre}       & \SI{0.560}{\metre} \\
Avg heading error       & \SI{4.6}{\degree}        & \SI{3.3}{\degree} \\
Time above \SI{5}{\metre\per\second} & 67\%        & 69\% \\
Wall collisions         & 0                        & 0 \\
Tests passing           & 6/6                      & 6/6 \\
\midrule
\multicolumn{3}{@{}l}{\textbf{FPGA Resources (Ultra96-V2 ZU3EG @ \SI{100}{\mega\hertz})}} \\
\midrule
DSP48E2                 & ---  & 321/360 (89.2\%) \\
LUT                     & ---  & 46,934/70,560 (66.5\%) \\
CLB                     & ---  & 8,491/8,820 (96.3\%) \\
FF                      & ---  & 39,081/141,120 (27.7\%) \\
Vivado WNS              & ---  & \SI{+0.044}{\nano\second} (timing closed) \\
\bottomrule
\end{tabular}
\end{table}

The FPGA achieves sub-\SI{1}{\milli\second} worst-case latency with deterministic timing (no OS jitter, cache misses, or thread preemption). The CPU MPC is faster per-call due to native float32 on x86, but the FPGA provides guaranteed real-time performance. Both implementations achieve zero wall collisions on the cl050 raceline.

The FPGA's higher lateral error (\SI{0.261}{\metre} vs \SI{0.110}{\metre}) is due to Q16.16 fixed-point quantization effects and the 19-step horizon (vs 20 for CPU). The heading error is actually lower on FPGA (\SI{3.3}{\degree} vs \SI{4.6}{\degree}).

\subsection{Architecture Notes}
\begin{itemize}
  \item \textbf{Communication bottleneck}: FPGA compute (\SI{0.828}{\milli\second}) is not the system bottleneck. The $\sim$\SI{2}{\milli\second} Jetson$\to$Ultra96 ROS2 network latency dominates end-to-end control latency.
  \item \textbf{Interrupt support}: The HLS IP includes interrupt hardware (GIE/IER/ISR registers + interrupt pin) for event-driven operation via UIO driver on Linux.
  \item \textbf{AXI-Lite interface}: Adequate for the 36-byte input payload at \SI{200}{\hertz}.
  \item \textbf{Clock limit}: \SI{100}{\mega\hertz} is dictated by the critical path (WNS = \SI{+0.044}{\nano\second}). Higher clocks would violate timing.
\end{itemize}

\end{document}
